\documentclass{article}
\usepackage{tekvideoexercises}

\begin{document}
\exercisename{Omskrivning: Negative eksponenter til brøker}
\streaklength{5}

\tableofcontents
\newpage

\begin{exercise}{Negative eksponenter 1-1}

Skriv $2^{-3}$ som en brøk.

\answerbox{\frac{1}{2^3}}

\hint

En negativ eksponent betyder den reciprokke (inverse) af den positive eksponent.

\hint

Regel: $a^{-n} = \frac{1}{a^n}$

\hint

Anvend reglen:

\hint
\[
2^{-3} = \frac{1}{2^3}
\]

\end{exercise}

\newpage

\begin{exercise}{Negative eksponenter 1-2}

Skriv $5^{-2}$ som en brøk.

\answerbox{\frac{1}{5^2}}

\hint

Anvend omskrivningsreglen:

\hint
\[
5^{-2} = \frac{1}{5^2}
\]

\end{exercise}

\newpage

\begin{exercise}{Negative eksponenter 1-3}

Skriv $x^{-4}$ som en brøk.

\answerbox{\frac{1}{x^4}}

\hint

Anvend reglen direkt:

\hint
\[
x^{-4} = \frac{1}{x^4}
\]

\end{exercise}

\newpage

\begin{exercise}{Negative eksponenter 1-4}

Skriv $3^{-1}$ som en brøk.

\answerbox{\frac{1}{3}}

\hint

Anvend reglen og reducér:

\hint
\[
3^{-1} = \frac{1}{3^1} = \frac{1}{3}
\]

\end{exercise}

\newpage

\begin{exercise}{Negative eksponenter 1-5}

Skriv $(a + b)^{-2}$ som en brøk.

\answerbox{\frac{1}{(a + b)^2}}

\hint

Anvend reglen for negative eksponenter:

\hint
\[
(a + b)^{-2} = \frac{1}{(a + b)^2}
\]

\end{exercise}

\newpage

\begin{exercise}{Negative eksponenter 1-6}

Skriv $4^{-3}$ som en brøk og beregn dens værdi.

\answerbox{\frac{1}{64}}

\hint

Omskriv til brøk:

\hint
\[
4^{-3} = \frac{1}{4^3}
\]

\hint

Beregn nævneren:

\hint
\[
4^3 = 64
\]

\hint

Færdiggør:

\hint
\[
\frac{1}{4^3} = \frac{1}{64}
\]

\end{exercise}

\end{document}
